\uuid{jxyB}
\titre{Application linéaire}
\auteur{hanani ; mzari}
\organisation{exo7}
\licence{CC BY-NC-SA 4.0}
\datecreate{2020-01-20}
\niveau{L1}
\module{Algèbre}
\chapitre{Application linéaire}
\sousChapitre{Définition}
\difficulte{3}
\theme{Applications linéaires, Application linéaire}
\exo7qcmid{lille-2.applications-lineaires.011}

\contenu{
\begin{question}
\qclassification{107.01}
\qtags{motcle=Application linéaire}
Quelles sont les assertions vraies ?
\begin{answers}
\good{Une application $f:\Rr \to \Rr$ est linéaire si et seulement s'il existe un réel $a$ tel que $f(x)=ax$, pour tout $x\in \Rr$}.
\bad{Une application $f:\Rr^2 \to \Rr^2$ est linéaire si et seulement s'il existe des réels $a$ et $b$ tels que $f(x,y)=(ax,by)$, pour tout $(x,y) \in \Rr^2$}.
\good{Une application $f:\Rr^2 \to \Rr^2$ est linéaire si et seulement s'il existe des réels $a,b,c$ et $d$ tels que $f(x,y)=(ax+by,cx+dy)$, pour tout $(x,y)\in \Rr^2$}.
\bad{Une application $f:\Rr^3 \to \Rr^3$ est linéaire si et seulement s'il existe des réels $a,b$ et $c$ tels que $f(x,y,z)=(ax,by,cz)$, pour tout $(x,y,z)\in \Rr^3$}.
\end{answers}
\begin{explanations} Une application $f:\Rr^n \to \Rr^m$ est linéaire si et seulement s'il existe des réels $a_{i,j}, 1\le i \le m, 1\le j \le n$,  tels que :\\ $f(x_1,x_2,\dots,x_n)=(a_{1,1}x_1+a_{1,2}x_2 + \dots + a_{1,n}x_n,\dots, a_{m,1}x_1+a_{m,2}x_2 + \dots + a_{m,n}x_n)$, pour tout $(x_1,x_2,\dots,x_n) \in \Rr^n$.
\end{explanations}
\end{question}
}
